% Provisional exposition; derived from the returned Pro proof and independently audited.
% Review/provenance status and exact integration instructions are in README.md.
\section{Independent linear deformations of the pentagon factors}
\label{sec:pentagon-affine}

The rotation formula leaves open what happens when the two pentagons are
stretched or sheared independently. Such changes generally alter both capacity
and volume. We show that they cannot increase the systolic ratio beyond the
Haim--Kislev--Ostrover value. Translations do not affect either quantity, so
linear maps suffice to describe the extension to affinely regular pentagons.

We use a regular pentagon with circumradius one and outward facet normals
\[
 n_i=(\cos(2\pi i/5),\sin(2\pi i/5)),\qquad
 P=\{q:n_i\cdot q\le h,\ 0\le i<5\},\qquad h=\cos(\pi/5).
\]
This normal convention differs by a common rotation from the preceding
rotation proof; applying that rotation to both coordinate planes is
symplectic. Its area is $S=\frac52\sin(2\pi/5)$. Throughout this section the product is
Lagrangian, with form
$\omega_0((q,p),(q',p'))=q\cdot p'-p\cdot q'$.

\begin{theorem}\label{thm:pentagon-affine-capacity}
Let $G,H\in GL(2,\mathbb R)$, and set
\[
 M=G^{-1}H^{-T},\qquad m(M)=\max_{0\le i,j<5}|n_i^TMn_j|.
\]
Then
\[
 c_{\mathrm{EHZ}}(GP\times HP)=\frac{(1+h)^2}{m(M)}.
\]
Moreover,
\[
 \operatorname{sys}(GP\times HP)
 :=\frac{c_{\mathrm{EHZ}}(GP\times HP)^2}
 {2\operatorname{area}(GP)\operatorname{area}(HP)}
 \le\frac{3+\sqrt5}{5}.
\]
Write $D=\operatorname{conv}\{\pm n_i:0\le i<5\}$ and $\alpha=\pi/10$.
Equality holds precisely when $M=tR_\alpha U$, where $t>0$ and $U$ is an
orthogonal symmetry of the regular decagon $D$.
\end{theorem}

The formula requires only twenty-five absolute pairings, even though the
capacity formula initially ranges over all facet orders and feasible weights.
The reason is a containment between convex sets of matrices. Unlike the
endpoint comparison for rotations, this containment is preserved by independent
linear changes in the two factors.

\subsection{Reducing the possible facet orders}
Write a polygon as $\{q:u_i\cdot q\le1\}$. Its normalized closure polytope is
\[
 C=\{\gamma\ge0:\ \textstyle\sum_i\gamma_i=1,\
                         \sum_i\gamma_i u_i=0\}.
\]
For the pentagon, $u_i=n_i/h$. The Haim--Kislev formula
\cite[Theorem 1.1]{HK2017}, with earlier word entries as the first argument of
$\omega_0$, reads
\[
 c_{\mathrm{EHZ}}=\frac1{2Q_{\max}},\qquad
 Q=\sum_{j<k}\beta_j\beta_k\omega_0(a_{\sigma_j},a_{\sigma_k}),
\]
where the product covectors $a_i$ are support-normalized, the weights have
sum one, and their weighted covector sum vanishes.

For a product, same-factor pairings vanish. If the total weight in the first
factor is $t$, the objective therefore has the form
$t(1-t)B(\gamma,\eta)$, where $B$ is bilinear in the normalized closure
weights of the factors. At a positive global maximum, maximizing successively
in $\gamma$ and $\eta$ allows both to be chosen at vertices of their closure
polytopes. The global value cannot increase, so these replacements preserve it.
A vertex has at most three positive entries: more than three columns of the
three closure equations admit a two-sided feasible perturbation. The remaining
factor $t(1-t)$ is maximal at $t=1/2$.

Thus some maximizer uses at most three facets of each factor. This is an
existence statement; it does not assert that all maximizing words are sparse.
Moving the first weighted covector cyclically to the end preserves $Q$, since
the change is its pairing with minus twice the sum of all the other weighted
covectors, and closure makes that pairing zero. We may consequently group
adjacent entries from the same factor cyclically. There are two, four, or six
alternating groups; the two-group value is zero.

Put $x_i=\gamma_i u_i$ and $y_j=\eta_j v_j$. These are factor steps with
sum zero; their actual product weights carry an additional factor $1/2$.
A four-group order $x,y,-x,-y$ has value $x\cdot y/2$. A six-group order
$x_1,y_1,x_2,y_2,x_3,y_3$ has value
\[
 Q=\frac12(x_1\cdot y_1-x_3\cdot y_2).
\]
The latter identity follows by pairing each $p$-step with its preceding
$q$-prefix: those prefixes are $x_1,-x_3,0$.
Represent these expressions by matrices
\[
 \frac12xy^T,\qquad \frac12(x_1y_1^T-x_3y_2^T),
\]
whose Frobenius pairing with the identity matrix is the objective. Under
independent changes $G,H$, both types transform by the same linear map
$T\mapsto G^{-T}TH^{-1}$. It remains to show that every six-group matrix for
regular pentagons lies in the convex hull of the four-group matrices.

\subsection{Two identities cover all six-group competitors}
Set
\[
 \lambda=\frac{\sqrt5-1}{2}=\frac1{2h},\qquad
 s=\frac1{1+h},\qquad \rho=\frac{s^2}{2},\qquad
 \mathcal D=\operatorname{conv}\{\pm n_i n_j^T:0\le i,j<5\}.
\]
A positive closure triple has all cyclic angular gaps less than $\pi$.
Its three positive integer index gaps therefore sum to five and are at most
two, so they are cyclic permutations of $(1,2,2)$. There are exactly five
triples, all rotations of $\{n_0,n_2,n_3\}$, and no opposite pairs.
The normalized steps for this representative triple are
\[
 se,sa,sb,\qquad e=n_0,\quad a=\lambda n_2,\quad
 b=\lambda n_3,\qquad e+a+b=0.
\]
Indeed the closure weights are $h/(1+h),1/[2(1+h)],1/[2(1+h)]$.
Every proper grouped sum of a triple is one step or its negative. The convex
hull of all such sums, over the five triples, is $sD$: all $\pm sn_i$ occur,
and the remaining steps are their contractions by $\lambda<1$.
It follows, by expanding independent convex combinations in each factor,
that the convex hull of the four-group matrices is exactly $\rho\mathcal D$.
Every generator is feasible: choose a triple with the specified singled-out
step, and group its other two steps together.

\begin{lemma}\label{lem:pentagon-affine-identities}
The two matrices $T_1=ee^T-aa^T$ and $T_2=eb^T-ae^T$ belong to $\mathcal D$.
\end{lemma}
\begin{proof}
They have the identities
\begin{align*}
 T_1&=\lambda^2ee^T-\lambda^2en_3^T-\lambda^3n_1n_2^T,\\
 T_2&=\lambda^2n_1n_3^T-\lambda^2n_2n_4^T.
\end{align*}
To verify them, substitute
$n_2+n_3=-\lambda^{-1}e$,
$\lambda n_1=e+n_2$, and $\lambda n_4=e+n_3$.
The first right-hand side becomes
$(\lambda^2+\lambda)ee^T-\lambda^2n_2n_2^T=T_1$;
the second becomes $\lambda(en_3^T-n_2e^T)=T_2$.
The unsigned coefficients in the first identity sum to
$2\lambda^2+\lambda^3=1$. Those in the second sum to
$2\lambda^2<1$, with the remaining mass assigned to $0\in\mathcal D$.
The signs are part of the defining generators, so these are convex
combinations, not signed combinations with negative weights.
\end{proof}

\begin{lemma}\label{lem:pentagon-affine-containment}
Every six-group matrix belongs to the four-group convex hull $\rho\mathcal D$.
\end{lemma}
\begin{proof}
Rotate each triple to the representative above and cut the word at the
$q$-step $se$. After division by $\rho$, the matrix is
$ev^T-wz^T$, where $w\in\{a,b\}$ and $v,z$ are distinct members of
$\{e,a,b\}$. Reflection $F=\operatorname{diag}(1,-1)$ on the left exchanges
$a,b$ and fixes $e$, so it suffices to take $w=a$. Its six possibilities are
\[
\begin{array}{c|cccccc}
(v,z)&(e,a)&(e,b)&(a,e)&(b,e)&(a,b)&(b,a)\\\hline
 ev^T-az^T&T_1&T_1F&T_2F&T_2&-FT_1F&-FT_1
\end{array}
\]
The last two equalities also follow from $e+a+b=0$. The set $\mathcal D$ is
invariant under sign, independent left and right reflections, and independent
pentagon rotations. Lemma~\ref{lem:pentagon-affine-identities} therefore
covers all twelve cyclic orders for every pair of closure triples.
\end{proof}

\subsection{Capacity and the sharp ratio}
\begin{proof}[Proof of Theorem~\ref{thm:pentagon-affine-capacity}]
The matrix containment survives the map $T\mapsto G^{-T}TH^{-1}$.
Thus the maximum is attained among transformed four-group matrices, giving
\[
 Q_{\max}=\rho\max_{i,j}|n_i^TG^{-1}H^{-T}n_j|.
\]
Taking $1/(2Q_{\max})$ proves the capacity formula. Six-group words may tie
this value; the argument only excludes a strictly larger value.

Since areas multiply by absolute determinants, the systolic ratio becomes
\[
 \frac{(1+h)^4}{2S^2}\frac{|\det M|}{m(M)^2}.
\]
Put $L=M/m(M)$. The inequalities defining $m(M)$ say exactly that
$LD\subseteq D^\circ$. The polar regular decagon is
$D^\circ=\sec\alpha\,R_\alpha D$, hence comparison of areas yields
\[
 |\det L|\le\sec^2\alpha.
\]
Now
\[
 \frac{(1+h)^4}{2S^2}=\frac{5+2\sqrt5}{10},\qquad
 \frac{5+2\sqrt5}{10}\sec^2(\pi/10)=\frac{3+\sqrt5}{5},
\]
which proves the bound. Equality of areas under convex-body containment
holds precisely when $LD=D^\circ$. Then
$\cos\alpha R_{-\alpha}L$ is a linear automorphism of $D$.
Such an automorphism permutes vertices and preserves or reverses cyclic
adjacency. Its action on two adjacent vertices determines it and agrees with
an orthogonal dihedral symmetry. This proves the stated equality
classification; conversely each listed matrix attains equality.
\end{proof}

For $G=I$ and $H=R_\theta$, the formula reduces to the rotation profile:
$m(R_\theta)=\cos d_5(\theta)$, where
$d_5(\theta)=\min_{k\in\mathbb Z}|\theta-k\pi/5|$.
The shorter endpoint proof remains useful because it establishes that profile
without the matrix classification. The original exhaustive SageMath argument
is retained as computational provenance in the repository's
\texttt{experiments/regular-products/} directory, under
\texttt{pentagon-rotation-formula-proof/}; it is not an input to either
analytic argument presented here.

The affine theorem concerns all independent linear images of regular
pentagons, including arbitrarily large deformations. The separate local
maximality theorem concerns nearby ten-facet polytopes and includes changes
that destroy the product structure. The two results control different sets of
perturbations. Neither is a bound for all ten-facet polytopes or for products
of arbitrary pentagons.
